\chapter{پیکربندی و قالب‌بندی خروجی}

در این فصل، بررسی ابزارهای متنی را ادامه داده و بر برنامه‌هایی تمرکز می‌کنیم که وظیفه‌ی آن‌ها آرایش و قالب‌بندی خروجی (\en{Output Formatting}) است، نه ایجاد تغییر در محتوای متن. این ابزارها عمدتاً جهت آماده‌سازی متون برای چاپ نهایی مورد استفاده قرار می‌گیرند؛ موضوعی که در فصل آتی به‌تفصیل به آن خواهیم پرداخت. ابزارهای مورد بررسی در این فصل عبارت‌اند از:

\begin{itemize}
  \item \cmd{nl}: شماره‌گذاری سطرهای یک فایل متنی (\en{Number Lines}).
  \item \cmd{fold}: شکستن سطرهای طولانی به خطوطی با طول مشخص، به‌منظور تناسب با عرض صفحه یا ترمینال.
  \item \cmd{fmt}: یک شکل‌دهنده‌ی متن (\en{Text Formatter}) ساده و سریع، برای بازآرایی و یکدست‌سازی پاراگراف‌ها.
  \item \cmd{pr}: آماده‌سازی متون برای عملیات چاپ، از طریق افزودن سرآیند، پانویس و صفحه‌بندی مناسب.
  \item \cmd{printf}: قالب‌بندی و چاپ داده‌ها بر اساس یک الگوی از پیش تعریف‌شده، با کنترل دقیق بر نحوه‌ی نمایش خروجی.
  \item \cmd{groff}: یک سیستم پیشرفته برای قالب‌بندی اسناد فنی و تولید خروجی چاپی با کیفیت بالا.
\end{itemize}

\section{ابزارهای پایه قالب‌بندی}

در ابتدا به بررسی برخی از ابزارهای ساده‌ی قالب‌بندی می‌پردازیم. این برنامه‌ها اغلب تک‌منظوره و دارای ساختاری ابتدایی هستند، اما برای انجام وظایف کوچک و استفاده در پایپ‌لاین‌ها (\en{Pipelines}) و اسکریپت‌های سیستمی، بسیار کاربردی و اثربخش عمل می‌کنند.

\subsection{دستور nl: شماره‌گذاری سطور}

ابزار \cmd{nl} برنامه‌ای کلاسیک برای انجام وظیفه‌ای ساده اما حیاتی است: شماره‌گذاری سطور متنی. در ابتدایی‌ترین حالت، عملکرد این دستور شباهت بسیاری به \cmd{cat -n} دارد:

\begin{latin}
\begin{lstlisting}
[me@linuxbox ~]$ nl distros.txt | head
     1  SUSE     10.2    12/07/2006
     2  Fedora   10      11/25/2008
     3  SUSE     11.0    06/19/2008
     4  Ubuntu   8.04    04/24/2008
     5  Fedora   8       11/08/2007
     6  SUSE     10.3    10/04/2007
     7  Ubuntu   6.10    10/26/2006
     8  Fedora   7       05/31/2007
     9  Ubuntu   7.10    10/18/2007
    10  Ubuntu   7.04    04/19/2007
\end{lstlisting}
\end{latin}

مشابه \cmd{cat} ابزار \cmd{nl} نیز می‌تواند چندین فایل را به عنوان آرگومان ورودی بپذیرد یا داده‌ها را از ورودی استاندارد (\en{Standard Input}) دریافت کند. با این حال، \cmd{nl} دارای گزینه‌های متعددی است و از نوعی نشانه‌گذاری مقدماتی برای مدیریت الگوهای پیچیده‌تر شماره‌گذاری پشتیبانی می‌کند.

این برنامه از مفهومی تحت عنوان «صفحات منطقی» (\en{Logical Pages}) پشتیبانی می‌کند که امکان بازنشانی (\en{Reset}) شمارش سطور را فراهم می‌آورد. با استفاده از گزینه‌های موجود، می‌توان شماره‌ی آغازین را روی مقدار خاصی تنظیم کرد و قالب نمایش آن را تا حدودی تغییر داد. در ساختار \cmd{nl} یک صفحه‌ی منطقی به سه بخش مجزا تقسیم می‌شود: سرآیند (\en{Header})، بدنه (\en{Body}) و پانویس (\en{Footer}).

در هر یک از این بخش‌ها، می‌توان شماره‌گذاری را از نو آغاز کرد یا سبک متفاوتی از نمایش را اعمال نمود. ذکر این نکته ضروری است که در صورت ارائه‌ی چندین فایل به \cmd{nl} این ابزار تمامی آن‌ها را به عنوان یک جریان متنی واحد (\en{Text Stream}) پردازش می‌کند. بخش‌های مختلف متن از طریق نشانه‌گذاری‌های ویژه‌ای (\en{Special Markup}) تفکیک می‌شوند که در جدول زیر تشریح شده است:

\begin{table}[htbp]
\centering
\caption{توالی‌های گریز در دستور \cmd{nl}}
\label{tab:nl-escapes}
\begin{tabularx}{\textwidth}{l X}
\toprule
\textbf{نشانه‌گذاری} & \textbf{مفهوم فنی} \\
\midrule
\cmd{\textbackslash:\textbackslash:\textbackslash:} & نشانه‌ی آغاز بخش سرآیند (\en{Header}) در یک صفحه‌ی منطقی (\en{Logical Page}). \\
\addlinespace
\cmd{\textbackslash:\textbackslash:} & نشانه‌ی آغاز بخش بدنه (\en{Body}) در یک صفحه‌ی منطقی. \\
\addlinespace
\cmd{\textbackslash:} & نشانه‌ی آغاز بخش پانویس (\en{Footer}) در یک صفحه‌ی منطقی. \\
\bottomrule
\end{tabularx}
\end{table}

هر یک از این نشانه‌گذاری‌های کنترلی (\en{Control Markup}) باید به‌صورت مجزا در یک سطر اختصاصی قرار گیرند. ابزار \cmd{nl} پس از شناسایی و پردازش این عناصر جهت مدیریت سلسله‌مراتب شماره‌گذاری، آن‌ها را از جریان خروجی نهایی حذف می‌کند.

\begin{table}[htbp]
\centering
\caption{گزینه‌های متداول دستور \cmd{nl}}
\label{tab:nl-options}
\begin{tabularx}{\textwidth}{l X}
\toprule
\textbf{گزینه} & \textbf{معنا} \\
\midrule
\cmd{-b style} & تعیین سبک شماره‌گذاری در بخش بدنه (\en{Body}). مقادیر مجاز عبارت‌اند از: \cmd{a} برای شماره‌گذاری تمامی سطرها؛ \cmd{t} برای شماره‌گذاری منحصراً سطرهای غیرخالی (حالت پیش‌فرض)؛ \cmd{n} برای غیرفعال کردن کامل شماره‌گذاری؛ و \cmd{pregexp} برای شماره‌گذاری صرفاً سطرهایی که با یک عبارت باقاعده‌ی (\en{Regular Expression}) مشخص مطابقت دارند. \\
\addlinespace
\cmd{-f style} & تعیین سبک شماره‌گذاری در بخش پانویس (\en{Footer})، با همان مقادیر مجازِ گزینه‌ی \cmd{-b}. حالت پیش‌فرض \cmd{n} (بدون شماره‌گذاری) است. \\
\addlinespace
\cmd{-h style} & تعیین سبک شماره‌گذاری در بخش سرآیند (\en{Header})، با همان مقادیر مجازِ گزینه‌ی \cmd{-b}. حالت پیش‌فرض \cmd{n} (بدون شماره‌گذاری) است. \\
\addlinespace
\cmd{-i number} & تعیین گام افزایشی (\en{Increment}) شماره‌گذاری، یعنی مقداری که شماره‌ی سطر در هر بار افزایش می‌یابد. مقدار پیش‌فرض برابر با \cmd{1} است. \\
\addlinespace
\cmd{-n format} & تعیین قالب نمایش شماره‌ی سطر. مقادیر مجاز عبارت‌اند از: \cmd{ln} برای نمایش چپ‌چین و بدون صفرهای پیشین (\en{Left Justified})؛ \cmd{rn} برای نمایش راست‌چین و بدون صفرهای پیشین (\en{Right Justified}، حالت پیش‌فرض)؛ و \cmd{rz} برای نمایش راست‌چین همراه با صفرهای پیشین (\en{Leading Zeros}). \\
\addlinespace
\cmd{-p} & غیرفعال کردن بازنشانی خودکار (\en{Reset}) شماره‌ها؛ در حالت عادی، شمارش در آغاز هر صفحه‌ی منطقی جدید از نو آغاز می‌شود، اما این گزینه باعث می‌شود شمارش به‌صورت پیوسته در سراسر سند ادامه یابد. \\
\addlinespace
\cmd{-s string} & تعیین رشته‌ای که به‌عنوان جداکننده، میان شماره‌ی سطر و محتوای آن درج می‌شود. کاراکتر پیش‌فرض برای این جداکننده، \en{Tab} است. \\
\addlinespace
\cmd{-v number} & تعیین شماره‌ی آغازینِ سطرها در ابتدای هر صفحه‌ی منطقی. مقدار پیش‌فرض برابر با \cmd{1} است. \\
\addlinespace
\cmd{-w width} & تعیین عرض ستونِ اختصاص‌یافته به شماره‌ی سطر، بر حسب تعداد کاراکتر. مقدار پیش‌فرض برابر با \cmd{6} است. \\
\bottomrule
\end{tabularx}
\end{table}

اگرچه نیاز به شماره‌گذاری سطور ممکن است به ندرت پیش بیاید، اما ترکیب \cmd{nl} با سایر ابزارهای متنی برای اجرای عملیات پیچیده بسیار سودمند است. در اینجا، گزارش توزیع‌های لینوکس را با بهره‌گیری از قابلیت «صفحات منطقی» (\en{Logical Pages}) تولید می‌کنیم. برای این منظور، اسکریپت \cmd{sed} فصل گذشته را به‌روزرسانی کرده و با نام \file{distros-nl.sed} ذخیره می‌کنیم:

\begin{latin}
\begin{lstlisting}
# sed script to produce Linux distributions report

1 i\
\\:\\:\\:\
\
Linux Distributions Report\
\
Name               Ver. Released\
----              ----- --------\
\\:\\:
s/\([0-9]\{2\}\)\/\([0-9]\{2\}\)\/\([0-9]\{4\}\)$/\3-\1-\2/
$ a\
\\:\
\
End Of Report
\end{lstlisting}
\end{latin}

این اسکریپت نشانه‌گذاری‌های مورد نیاز \cmd{nl} برای تفکیک صفحات منطقی را درج کرده و یک پانویس (\en{Footer}) نیز به انتهای گزارش می‌افزاید. توجه داشته باشید که برای جلوگیری از تفسیر نادرست نویسه‌ی بک‌اسلش (\cmd{\textbackslash}) توسط \cmd{sed} از حالت گریز دوگانه (\cmd{\textbackslash\textbackslash}) استفاده شده است.

در ادامه، با ایجاد یک پایپ‌لاین (\en{Pipeline}) متشکل از \cmd{sort}، \cmd{sed} و \cmd{nl} گزارش نهایی را استخراج می‌کنیم:

\begin{latin}
\begin{lstlisting}
[me@linuxbox ~]$ sort -k 1,1 -k 2n distros.txt | sed -f distros-nl.sed | nl

     Linux Distributions Report

       Name    Ver.    Released
       ----    ----    --------

     1  Fedora  5       2006-03-20
     2  Fedora  6       2006-10-24
     3  Fedora  7       2007-05-31
     4  Fedora  8       2007-11-08
     5  Fedora  9       2008-05-13
     6  Fedora  10      2008-11-25
     7  SUSE    10.1    2006-05-11
     8  SUSE    10.2    2006-12-07
     9  SUSE    10.3    2007-10-04
    10  SUSE    11.0    2008-06-19
    11  Ubuntu  6.06    2006-06-01
    12  Ubuntu  6.10    2006-10-26
    13  Ubuntu  7.04    2007-04-19
    14  Ubuntu  7.10    2007-10-18
    15  Ubuntu  8.04    2008-04-24
    16  Ubuntu  8.10    2008-10-30


     End Of Report
\end{lstlisting}
\end{latin}

ساختار این گزارش نتیجه‌ی یک پردازش زنجیره‌ای است: ابتدا فهرست توزیع‌ها بر اساس نام و نسخه (فیلدهای ۱ و ۲) توسط \cmd{sort} مرتب می‌شود؛ سپس \cmd{sed} سرآیند، پانویس و نشانه‌گذاری‌های منطقی را به جریان متن تزریق می‌کند. در نهایت، \cmd{nl} خروجی را دریافت کرده و طبق رفتار پیش‌فرض خود، منحصراً سطور بخش بدنه (\en{Body}) را شماره‌گذاری می‌کند.

شما می‌توانید این فرمان را با استفاده از گزینه‌های متنوع \cmd{nl} جهت شخصی‌سازی قالب خروجی آزمایش کنید. برای نمونه، استفاده از دستور:

\begin{latin}
\begin{lstlisting}
nl -n rz
\end{lstlisting}
\end{latin}

باعث می‌شود شماره‌گذاری‌ها به‌صورت راست‌چین و همراه با صفرهای پیشین (\en{Leading Zeros}) نمایش داده شوند. همچنین با به‌کارگیری دستور:

\begin{latin}
\begin{lstlisting}
nl -w 3 -s ' '
\end{lstlisting}
\end{latin}

می‌توانید عرض فیلد شماره‌گذاری را به ۳ کاراکتر محدود کرده و از یک کاراکتر فاصله (\en{Space}) به عنوان جداکننده میان شماره و متن استفاده کنید.

\subsection{دستور fold: شکستن سطور}

دستور \cmd{fold}، که معنای واژگانی آن «تا کردن» است، وظیفه‌ی شکستن سطرهای متن را برای انطباق با یک عرض (\en{Width}) مشخص بر عهده دارد. این ابزار نیز، همانند سایر دستورات مشابه، داده‌های ورودی را از یک یا چند فایل متنی، یا از طریق ورودی استاندارد (\en{Standard Input})، دریافت می‌کند. برای درک بهتر سازوکار این دستور، یک جریان متنی ساده را به آن ارسال می‌کنیم:

\begin{latin}
\begin{lstlisting}
[me@linuxbox ~]$ echo "The quick brown fox jumps over the lazy dog." | fold -w 12
The quick br
own fox jump
s over the l
azy dog.
\end{lstlisting}
\end{latin}

در این مثال، متن ارسالی توسط دستور \cmd{echo} به قطعاتی تقسیم شده است که عرض هر یک توسط گزینه‌ی \cmd{-w} تعیین می‌گردد. در اینجا، عرض سطر را ۱۲ نویسه (\en{Character}) در نظر گرفته‌ایم. چنانچه مقدار عرض مشخص نشود، \cmd{fold} به‌صورت پیش‌فرض از استاندارد ۸۰ نویسه استفاده می‌کند.

نکته‌ی قابل توجه این است که در حالت پیش‌فرض، شکستن سطور بدون در نظر گرفتن مرز کلمات انجام می‌شود. برای رفع این مشکل، افزودن گزینه‌ی \cmd{-s} به \cmd{fold} فرمان می‌دهد که سطر را در آخرین فضای خالی (\en{Space}) موجود، پیش از رسیدن به حد نصاب عرض تعیین‌شده، بشکند تا یکپارچگی کلمات حفظ شود:

\begin{latin}
\begin{lstlisting}
[me@linuxbox ~]$ echo "The quick brown fox jumps over the lazy dog." | fold -w 12 -s
The quick
brown fox
jumps over
the lazy
dog.
\end{lstlisting}
\end{latin}

\subsection{دستور fmt: قالب‌بندی پاراگراف‌ها}

ابزار \cmd{fmt} همانند \cmd{fold} وظیفه‌ی شکستن سطور را بر عهده دارد، اما قابلیت‌های پیشرفته‌تری را برای مدیریت متن ارائه می‌دهد. این برنامه با دریافت فایل یا ورودی استاندارد (\en{Standard Input})، عملیات آرایش پاراگراف‌ها را به شکلی پویا انجام می‌دهد؛ به این صورت که با حفظ خطوط خالی و تورفتگی‌ها (\en{Indentation})، فواصل میان واژگان را پر کرده و سطور را به هم متصل می‌کند.

برای بررسی عملکرد این ابزار، از بخشی از مستندات رسمی (\cmd{info fmt}) استفاده می‌کنیم:

\begin{latin}
\begin{lstlisting}
   `fmt` reads from the specified FILE arguments (or standard input if none are given), and writes to standard output.

    By default, blank lines, spaces between words, and indentation are preserved in the output; successive input lines with different indentation are not joined; tabs are expanded on input and introduced on output.

   `fmt` prefers breaking lines at the end of a sentence, and tries to avoid line breaks after the first word of a sentence or before the last word of a sentence. A "sentence break" is defined as either the end of a paragraph or a word ending in any of `.?!', followed by two spaces or end of line, ignoring any intervening parentheses or quotes. Like TeX, `fmt` reads entire "paragraphs" before choosing line breaks; the algorithm is a variant of that given by Donald E. Knuth and Michael F. Plass in "Breaking Paragraphs Into Lines", `Software--Practice & Experience` 11, 11 (November 1981), 1119-1184.
\end{lstlisting}
\end{latin}

این متن را در یک فایل با نام \file{fmt-info.txt} ذخیره می‌کنیم. فرض کنید قصد داریم این محتوا را به‌گونه‌ای قالب‌بندی کنیم که در ستونی با عرض حداکثری ۵۰ نویسه (\en{Character}) قرار گیرد. برای این منظور، از ابزار \cmd{fmt} به همراه گزینه \cmd{-w} استفاده می‌کنیم:

\begin{latin}
\begin{lstlisting}
[me@linuxbox ~]$ fmt -w 50 fmt-info.txt | head
    `fmt' reads from the specified FILE arguments
(or standard input if
none are given), and writes to standard output.

    By default, blank lines, spaces between words,
and indentation are
preserved in the output; successive input lines
with different indentation are not joined; tabs
are expanded on input and introduced on output.
\end{lstlisting}
\end{latin}

در نگاه اول، خروجی حاصل ممکن است کمی غیرمنتظره به نظر برسد. برای درک علت این نوع نمایش، بازخوانی دقیق متن راهنما ضروری است، چرا که دقیقاً منطقِ زیربنایی این نوع پردازش را تبیین می‌کند:

\begin{quote}
«به‌صورت پیش‌فرض، خطوط خالی، فواصل میان کلمات و تورفتگی‌ها در خروجی حفظ می‌شوند؛ سطور ورودی متوالی که دارای تورفتگی‌های متفاوتی هستند، به یکدیگر متصل نمی‌شوند؛ نویسه‌های \en{Tab} در ورودی بسط داده شده و در خروجی درج می‌گردند.»
\end{quote}

بنابراین، ابزار \cmd{fmt} در حال حفظ تورفتگی سطر نخست است. خوشبختانه، این برنامه گزینه‌ای برای اصلاح این رفتار و مدیریت بهینه‌ی حاشیه‌ها در اختیار قرار می‌دهد:

\begin{latin}
\begin{lstlisting}
[me@linuxbox ~]$ fmt -cw 50 fmt-info.txt
    `fmt' reads from the specified FILE arguments
(or standard input if none are given), and writes
to standard output.

    By default, blank lines, spaces between words,
and indentation are preserved in the output;
successive input lines with different indentation
are not joined; tabs are expanded on input and
introduced on output.

    `fmt' prefers breaking lines at the end of a
sentence, and tries to avoid line breaks after
the first word of a sentence or before the last
word of a sentence. A "sentence break" is defined
as either the end of a paragraph or a word ending
in any of `.?!', followed by two spaces or end of
line, ignoring any intervening parentheses or
quotes. Like TeX, `fmt' reads entire "paragraphs"
before choosing line breaks; the algorithm is a
variant of that given by Donald E. Knuth and
Michael F. Plass in "Breaking Paragraphs Into
Lines", `Software--Practice & Experience' 11, 11
(November 1981), 1119-1184.
\end{lstlisting}
\end{latin}

این نتیجه بسیار مطلوب‌تر است. با افزودن گزینه \cmd{-c} (مخفف \en{Crown Margin})، اکنون خروجی دقیقاً مطابق با ساختار مورد نظر قالب‌بندی شده است. در این حالت، \cmd{fmt} تورفتگی سطر اول پاراگراف را تشخیص داده و سایر سطور را بر اساس آن تراز می‌کند.

ابزار \cmd{fmt} گزینه‌های کاربردی و منعطفی را ارائه می‌دهد که در جدول زیر به تشریح آن‌ها می‌پردازیم:

\begin{table}[htbp]
\centering
\caption{گزینه‌های عملیاتی دستور \cmd{fmt}}
\label{tab:fmt-options}
\begin{tabularx}{\textwidth}{l X}
\toprule
\textbf{گزینه} & \textbf{شرح عملکرد تخصصی} \\
\midrule
\cmd{-c} & فعال کردن حالت \en{Crown Margin}؛ در این حالت، تورفتگی (\en{Indentation}) دو سطر نخست هر پاراگراف حفظ می‌شود و سطرهای بعدی آن پاراگراف با میزان تورفتگیِ سطر دوم هم‌تراز می‌گردند. \\
\addlinespace
\cmd{-p string} & قالب‌بندی بر اساس پیشوند (\en{Prefix})؛ در این حالت، تنها سطرهایی که با رشته‌ی متنی مشخص‌شده (\file{string}) آغاز می‌شوند، مشمول عملیات قالب‌بندی قرار می‌گیرند و سایر سطرها بدون تغییر باقی می‌مانند. پس از پایان پردازش، همان پیشوند مجدداً به ابتدای هر سطر افزوده می‌شود. این قابلیت، به‌طور خاص، برای قالب‌بندی کامنت‌ها در کدهای برنامه‌نویسی یا فایل‌های پیکربندی (برای نمونه، سطرهایی با پیشوند \cmd{\#}) بسیار کاربردی است. \\
\addlinespace
\cmd{-s} & فعال کردن حالت شکست سطر تنها (\en{Split Only})؛ در این حالت، سطرهای طولانی صرفاً برای انطباق با عرض تعیین‌شده شکسته می‌شوند، بدون آنکه سطرهای کوتاه‌تر برای پر کردن فضای خالی با یکدیگر ترکیب شوند. این گزینه به‌ویژه برای قالب‌بندی کدهای برنامه‌نویسی، که حفظ ساختار اصلی خطوط در آن‌ها ضروری است، کاربرد دارد. \\
\addlinespace
\cmd{-u} & فعال کردن فاصله‌گذاری یکنواخت (\en{Uniform Spacing})؛ این گزینه استاندارد کلاسیک ماشین‌تحریر را اعمال می‌کند، به این صورت که میان کلمات یک فاصله و میان جملات دو فاصله قرار می‌گیرد. این قابلیت برای رفع ناهمواری‌های ناشی از تراز اجباری متن (\en{Justification}) و حذف فاصله‌های نامتعارف مفید است. \\
\addlinespace
\cmd{-w width} & تعیین عرض ستون؛ این گزینه متن را برای انطباق با عرض مشخصی از نویسه‌ها قالب‌بندی می‌کند. مقدار پیش‌فرض این عرض برابر با ۷۵ نویسه است. \\
\bottomrule
\end{tabularx}
\end{table}

\begin{note}
\textbf{نکته:} ابزار \cmd{fmt}، به‌منظور حفظ توازن دیداری خروجی و پیشگیری از بروز شکستگی‌های نامنظم در انتهای سطرها، معمولاً طول خطوط را اندکی کمتر از مقدار تعیین‌شده در گزینه‌ی \cmd{-w} تنظیم می‌کند.
\end{note}

گزینه \cmd{-p} قابلیتی هوشمند و بسیار کاربردی است. با بهره‌گیری از این گزینه، می‌توان بخش‌های معینی از یک فایل را به شرطی که با رشته‌ای (\en{Prefix}) یکسان آغاز شوند، قالب‌بندی مجدد کرد. از آنجا که در بسیاری از زبان‌های برنامه‌نویسی و فایل‌های پیکربندی، علامت \cmd{\#} نشان‌دهنده‌ی آغاز کدهای توضیحی (\en{Comments}) است، این گزینه ابزاری ایده‌آل برای مرتب‌سازی آن‌ها محسوب می‌شود.

برای بررسی این قابلیت، ابتدا فایلی نمونه ایجاد می‌کنیم:

\begin{latin}
\begin{lstlisting}
[me@linuxbox ~]$ cat > fmt-code.txt
# This file contains code with comments.

# This line is a comment.
# Followed by another comment line.
# And another.

This, on the other hand, is a line of code.
And another line of code.
And another.
\end{lstlisting}
\end{latin}

در این فایل، خطوط توضیحی با رشته‌ی \cmd{\#} (علامت هش به همراه یک فاصله) آغاز شده‌اند، در حالی که خطوط مربوط به بدنه اصلی کد فاقد این پیشوند هستند. اکنون با استفاده از \cmd{fmt} و گزینه \cmd{-p} تنها بخش‌های توضیحی را قالب‌بندی کرده و ساختار کد را دست‌نخورده باقی می‌گذاریم:

\begin{latin}
\begin{lstlisting}
[me@linuxbox ~]$ fmt -w 50 -p '# ' fmt-code.txt
# This file contains code with comments.

# This line is a comment. Followed by another
# comment line. And another.

This, on the other hand, is a line of code.
And another line of code.
And another.
\end{lstlisting}
\end{latin}

همان‌طور که مشاهده می‌کنید، سطور مربوط به توضیحات برای انطباق با عرض تعیین‌شده با یکدیگر ترکیب شده‌اند، در حالی که خطوط خالی و سطوری که فاقد پیشوند مشخص‌شده بودند، بدون هیچ تغییری در جای خود باقی مانده‌اند. این ویژگی به توسعه‌دهندگان اجازه می‌دهد تا مستندات داخل کد را بدون آسیب رساندن به منطق برنامه، به‌طور خودکار مرتب کنند.

\subsection{دستور pr: آماده‌سازی متن برای چاپ}

ابزار \cmd{pr} برای صفحه‌بندی (\en{Pagination}) متون طراحی شده است. هنگام ارسال یک متن به چاپگر، معمولاً ترجیح داده می‌شود که صفحات با مقداری فضای خالی (\en{Whitespace}) در بالا و پایین از یکدیگر متمایز شوند تا حاشیه‌های لازم ایجاد گردد. علاوه بر این، می‌توان از این فضاهای خالی برای درج سرآیند (\en{Header}) و پانویس (\en{Footer}) در هر صفحه بهره جست.

برای نمایش عملکرد \cmd{pr} از فایل \file{distros.txt} استفاده می‌کنیم و آن را در قالب مجموعه‌ای از صفحات کوتاه پیکربندی می‌نماییم (در اینجا تنها دو صفحه‌ی نخست نمایش داده شده است):

\begin{latin}
\begin{lstlisting}
[me@linuxbox ~]$ pr -l 15 -w 65 distros.txt

2025-12-11 18:27                distros.txt                  Page 1

SUSE    10.2      12/07/2006
Fedora  10        11/25/2008
SUSE    11.0      06/19/2008
Ubuntu  8.04      04/24/2008
Fedora  8         11/08/2007


2025-12-11 18:27                distros.txt                  Page 2

SUSE    10.3      10/04/2007
Ubuntu  6.10      10/26/2006
Fedora  7         05/31/2007
Ubuntu  7.10      10/18/2007
Ubuntu  7.04      04/19/2007
\end{lstlisting}
\end{latin}

در این نمونه، از گزینه \cmd{-l} (جهت تعیین طول صفحه) و گزینه \cmd{-w} (جهت تعیین عرض صفحه) استفاده شده است تا صفحاتی با عرض ۶۵ ستون و طول ۱۵ سطر تعریف شوند. ابزار \cmd{pr} محتوای فایل را صفحه‌بندی کرده، میان صفحات فواصل لازم را ایجاد می‌کند و یک سرآیند پیش‌فرض شامل زمان آخرین ویرایش فایل، نام فایل و شماره‌ی صفحه تولید می‌نماید.

این برنامه گزینه‌های متعددی برای کنترل دقیق چیدمان صفحه در اختیار کاربر قرار می‌دهد که در فصل ۲۲ به‌صورت تفصیلی آن‌ها را بررسی خواهیم کرد.

\subsection{دستور printf: قالب‌بندی و چاپ داده‌ها}

برخلاف دیگر دستورات این فصل، دستور \cmd{printf} برای استفاده در پایپ‌لاین (\en{pipeline}) طراحی نشده است (ورودی استاندارد را نمی‌پذیرد) و کاربرد رایجی در خط فرمان ندارد (اغلب در اسکریپت‌ها استفاده می‌شود). پس چرا مهم است؟ چون کاربرد بسیار گسترده‌ای دارد.

\cmd{printf} (از عبارت \en{``print formatted''}) در ابتدا برای زبان برنامه‌نویسی \en{C} توسعه یافت و در بسیاری از زبان‌های برنامه‌نویسی از جمله \cmd{shell} پیاده‌سازی شده است. در واقع، در \cmd{bash}، دستور \cmd{printf} به صورت داخلی (\en{builtin}) موجود است.

فرم کلی استفاده از \cmd{printf} به این صورت است:

\begin{latin}
\begin{lstlisting}
printf [-v var] "format" arguments
\end{lstlisting}
\end{latin}

در این ساختار، یک رشته‌ی کنترلی تحت عنوان «قالب» (\en{format}) تعریف می‌شود که قواعد آرایش داده‌ها را بر روی لیستی از آرگومان‌ها اعمال می‌کند. خروجی نهایی به‌صورت پیش‌فرض به خروجی استاندارد ارسال می‌گردد، مگر اینکه با استفاده از گزینه \cmd{-v} ذخیره‌سازی نتیجه در یک متغیر درخواست شود.

در ادامه، یک نمونه‌ی ساده از کاربرد این دستور ارائه شده است:

\begin{latin}
\begin{lstlisting}
[me@linuxbox ~]$ printf "I formatted the string: %s\n" foo
I formatted the string: foo
\end{lstlisting}
\end{latin}

در این مثال، \cmd{\%s} یک جای‌نگهدار (\en{Placeholder}) برای رشته است که آرگومان \cmd{foo} جایگزین آن می‌شود و \cmd{\textbackslash n} نشان‌دهنده‌ی کاراکتر خط جدید است.

در ادامه، نمونه‌ی قبل را با استفاده از گزینه \cmd{-v} ملاحظه می‌کنید؛ این گزینه به‌جای ارسال مستقیم خروجی به ترمینال، نتیجه‌ی پردازش‌شده را در یک متغیر ذخیره می‌کند:

\begin{latin}
\begin{lstlisting}
[me@linuxbox ~]$ printf -v result "I formatted the string: %s\n" foo
[me@linuxbox ~]$ echo "$result"
I formatted the string: foo
\end{lstlisting}
\end{latin}

رشته‌ی قالب (\en{Format String}) می‌تواند ترکیبی از متن ثابت (مانند \en{``I formatted the string''})، توالی‌های گریز (مانند \cmd{\textbackslash n} برای ایجاد سطر جدید) و مجموعه‌ای از کاراکترهای ویژه باشد که با نماد \cmd{\%} آغاز می‌شوند؛ این توالی‌ها در ادبیات برنامه‌نویسی «مشخصه‌های قالب‌بندی» (\en{Conversion Specifications}) نامیده می‌شوند. در مثال فوق، مشخصه‌ی \cmd{\%s} وظیفه‌ی قالب‌بندی رشته‌ی \cmd{"foo"} و جای‌گذاری آن در موقعیت مقرر را بر عهده دارد. نمونه‌ی دیگری از این سازوکار را در زیر مشاهده می‌کنید:

\begin{latin}
\begin{lstlisting}
[me@linuxbox ~]$ printf "I formatted '%s' as a string.\n" foo
I formatted 'foo' as a string.
\end{lstlisting}
\end{latin}

همان‌طور که مشهود است، مشخصه‌ی تبدیل \cmd{\%s} در خروجی نهایی با مقدار رشته‌ای \cmd{"foo"} جایگزین شده است. در حالی که کاراکتر \cmd{s} منحصراً برای آرایش داده‌های متنی (\en{String}) به کار می‌رود، مشخصه‌های متنوع دیگری نیز برای مدیریت و نمایش سایر گونه‌های داده (مانند اعداد صحیح یا ممیز شناور) تعریف شده‌اند. در جدول زیر، پرکاربردترین مشخصه‌های تعیین نوع داده (\en{Data Type}) در دستور \cmd{printf} فهرست شده است:

\begin{table}[htbp]
\centering
\caption{مشخصه‌های تعیین نوع داده در \cmd{printf}}
\label{tab:printf-types}
\begin{tabularx}{\textwidth}{c X}
\toprule
\textbf{مشخصه} & \textbf{شرح عملکرد فنی} \\
\midrule
\cmd{d} & قالب‌بندی و چاپ عدد به‌صورت یک عدد صحیح اعشاری علامت‌دار (\en{Signed Decimal Integer}). \\
\addlinespace
\cmd{f} & قالب‌بندی و چاپ عدد به‌صورت ممیز شناور (\en{Floating-point Number})، با شش رقم اعشار به‌عنوان دقت پیش‌فرض. \\
\addlinespace
\cmd{o} & تبدیل و قالب‌بندی عدد صحیح به مبنای ۸ (اُکتال). \\
\addlinespace
\cmd{s} & قالب‌بندی و چاپ داده‌های متنی (رشته). \\
\addlinespace
\cmd{x} & تبدیل و قالب‌بندی عدد صحیح به مبنای ۱۶ (هگزادسیمال)، با نمایش حروف به‌صورت کوچک (\en{a} تا \en{f}). \\
\addlinespace
\cmd{X} & مشابه \cmd{x}؛ با این تفاوت که حروف هگزادسیمال به‌صورت بزرگ (\en{A} تا \en{F}) نمایش داده می‌شوند. \\
\addlinespace
\cmd{\%} & چاپ خودِ نویسه‌ی \cmd{\%} به‌صورت متنی؛ برای این منظور باید از توالی \cmd{\%\%} در رشته‌ی قالب استفاده کرد. \\
\bottomrule
\end{tabularx}
\end{table}

در ادامه، تأثیر هر یک از این مشخصه‌های تبدیل بر مقدار عددی \cmd{380} مشاهده می‌شود:

\begin{latin}
\begin{lstlisting}
[me@linuxbox ~]$ printf "%d, %f, %o, %s, %x, %X\n" 380 380 380 380
380, 380.000000, 574, 380, 17c, 17C
\end{lstlisting}
\end{latin}

از آنجا که شش مشخصه‌ی تبدیل در رشته‌ی قالب تعریف شده‌اند، باید شش آرگومان متناظر نیز به \cmd{printf} ارائه شود. نتایج به‌دست‌آمده، تفاوت در نحوه‌ی نمایش یک مقدار واحد را بر اساس نوع مشخصه نشان می‌دهند.

برای کنترل دقیق‌تر خروجی، می‌توان مؤلفه‌های اختیاری متعددی را به مشخصه‌ی تبدیل افزود. ساختار کامل یک مشخصه‌ی تبدیل به ترتیب اولویت زیر است:

\begin{latin}
\begin{lstlisting}
%[flags][width][.precision]conversion_specification
\end{lstlisting}
\end{latin}

هنگام به‌کارگیری هم‌زمان این پارامترها، رعایت ترتیب فوق برای تفسیر صحیح توسط پوسته الزامی است. جدول زیر نقش هر یک از این اجزا را تبیین می‌کند:

\begin{table}[htbp]
\centering
\caption{اجزای اختیاری مشخصه‌های تبدیل \cmd{printf}}
\label{tab:printf-components}
\begin{tabularx}{\textwidth}{l X}
\toprule
\textbf{مؤلفه} & \textbf{توضیح} \\
\midrule
\en{flags} & شامل پنج پرچم عملیاتی برای تنظیم ظاهر خروجی:
\begin{itemize}
  \item \cmd{\#} (قالب جایگزین): نحوه نمایش را بر اساس نوع داده تغییر می‌دهد؛ در مبنای ۸ (\cmd{o}) یک صفر (\cmd{0}) و در مبنای ۱۶ (\cmd{x} یا \cmd{X}) پیشوندهای \cmd{0x} یا \cmd{0X} را به عدد اضافه می‌کند. (این پیشوند تنها زمانی درج می‌شود که مقدار عدد غیرصفر باشد).
  \item \cmd{0} (پرکننده‌ی صفر): فیلد خروجی را با صفرهای پیشین (\en{Leading Zeros}) پر می‌کند (مانند \cmd{000380}).
  \item \cmd{-} (تراز چپ): خروجی را به سمت چپ تراز می‌کند (حالت پیش‌فرض \cmd{printf} تراز راست است).
  \item \en{` `} (نویسه‌ی فاصله): برای اعداد مثبت، یک فضای خالی در محل علامت عدد قرار می‌دهد.
  \item \cmd{+} (نمایش علامت): باعث می‌شود اعداد مثبت نیز با علامت \cmd{+} نمایش داده شوند (در حالت عادی فقط اعداد منفی علامت‌دار هستند).
\end{itemize} \\
\addlinespace
\en{width} & عددی که حداقل عرض فیلد (\en{Field Width}) را بر حسب نویسه (\en{Character}) تعیین می‌کند. چنانچه طول واقعی خروجی از این مقدار کمتر باشد، فضای خالی باقی‌مانده با فاصله یا صفر (بسته به پرچم‌های اعمال‌شده) پر می‌شود. \\
\addlinespace
\cmd{.precision} & در اعداد ممیز شناور، تعداد ارقام اعشار قابل‌نمایش را مشخص می‌کند. در مورد رشته‌ها (\en{Strings})، این پارامتر تعیین‌کننده‌ی حداکثر تعداد نویسه‌هایی است که از رشته در خروجی نمایش داده می‌شود. \\
\bottomrule
\end{tabularx}
\end{table}

جدول زیر نمونه‌هایی از پیاده‌سازی مشخصه‌های تبدیل و تأثیر آن‌ها بر خروجی نهایی را نمایش می‌دهد:

\begin{table}[htbp]
\centering
\caption{نمونه‌های کاربردی مشخصه‌های تبدیل در \cmd{printf}}
\label{tab:printf-examples}
\begin{tabularx}{\textwidth}{l l l X}
\toprule
\textbf{آرگومان} & \textbf{قالب (\en{Format})} & \textbf{نتیجه} & \textbf{تحلیل فنی} \\
\midrule
\cmd{380} & \cmd{"\%d"} & \cmd{380} & نگاشت ساده‌ی مقدار به یک عدد صحیح ده‌دهی، بدون هیچ‌گونه پردازش اضافی. \\
\addlinespace
\cmd{380} & \cmd{"\%\#x"} & \cmd{0x17c} & تبدیل عدد به مبنای ۱۶ (هگزادسیمال)، همراه با درج خودکار پیشوند \cmd{0x} که ناشی از به‌کارگیری پرچم \cmd{\#} است. \\
\addlinespace
\cmd{380} & \cmd{"\%05d"} & \cmd{00380} & تعیین حداقل عرض ۵ نویسه برای فیلد، و پر کردن فضای خالی ابتدای عدد با صفر (\en{Padding}). \\
\addlinespace
\cmd{380} & \cmd{"\%05.5f"} & \cmd{380.00000} & نمایش عدد با دقت ۵ رقم اعشار؛ از آنجا که طول خروجی حاصل (۹ نویسه) از عرض فیلد تعیین‌شده (۵ نویسه) بیشتر است، هیچ فاصله یا کاراکتر اضافی برای پر کردن عرض فیلد درج نمی‌شود. \\
\addlinespace
\cmd{380} & \cmd{"\%010.5f"} & \cmd{0380.00000} & با افزایش عرض فیلد به ۱۰ نویسه، و از آنجا که طول خروجی (۹ نویسه) کمتر از این مقدار است، یک صفر پیشین برای تکمیل ظرفیت فیلد افزوده می‌شود. \\
\addlinespace
\cmd{380} & \cmd{"\%+d"} & \cmd{+380} & اجبار به نمایش علامت مثبت پیش از عدد، که در حالت عادی (بدون پرچم \cmd{+}) نمایش داده نمی‌شود. \\
\addlinespace
\cmd{380} & \cmd{"\%-d"} & \cmd{380} & تراز کردن مقدار به سمت چپ فیلد؛ تأثیر این پرچم تنها زمانی قابل مشاهده است که عرض فیلد از طول واقعی عدد بیشتر باشد؛ در این حالت، فاصله‌های خالی در سمت راست عدد اضافه می‌شوند. \\
\addlinespace
\cmd{abcdefghijk} & \cmd{"\%5s"} & \cmd{abcdefghijk} & تعیین حداقل عرض ۵ نویسه برای رشته؛ از آنجا که طول رشته‌ی ورودی بیش از این مقدار است، خروجی بدون هیچ تغییری چاپ می‌شود. \\
\addlinespace
\cmd{abcdefghijk} & \cmd{"\%.5s"} & \cmd{abcde} & استفاده از پارامتر دقت (\en{Precision}) برای بریدن (\en{Truncate}) رشته و محدود کردن آن به ۵ نویسه‌ی نخست. \\
\bottomrule
\end{tabularx}
\end{table}

مجدداً تأکید می‌شود که در دستور \cmd{printf}، عرض فیلد (\en{width}) تنها حداقل فضای اختصاص‌یافته به خروجی را تضمین می‌کند و هرگز باعث قطع‌شدن یا کوتاه‌شدن داده نمی‌شود؛ حتی اگر طول واقعی داده از عرض تعیین‌شده بیشتر باشد، کل داده بدون هیچ کاستی نمایش داده خواهد شد. در مقابل، پارامتر دقت (\en{precision}) رفتاری کاملاً متفاوت دارد: در مورد رشته‌های متنی، این پارامتر می‌تواند باعث بریدن (\en{Truncate}) بخشی از داده و حذف نویسه‌های اضافی شود؛ و در مورد اعداد اعشاری، تعداد ارقام پس از ممیز را کنترل کرده و در صورت لزوم، مقدار عدد را گرد (\en{Round}) می‌کند.

در اغلب موارد، \cmd{printf} برای قالب‌بندی داده‌های جدولی در اسکریپت‌ها به کار می‌رود و کمتر به‌صورت مستقیم در خط فرمان استفاده می‌شود. با این حال، می‌توان از آن برای حل چالش‌های مختلف چیدمان متن بهره برد. ابتدا نحوه‌ی جداسازی فیلدها با استفاده از نویسه‌ی \en{Tab} را بررسی می‌کنیم:

\begin{latin}
\begin{lstlisting}
[me@linuxbox ~]$ printf "%s\t%s\t%s\n" str1 str2 str3
str1	str2	str3
\end{lstlisting}
\end{latin}

با درج توالی گریز \cmd{\textbackslash t} (\en{Escape Sequence} برای \en{Tab})، خروجی ستون‌بندی می‌شود. در ادامه، نمونه‌ای از قالب‌بندی منظم اعداد را مشاهده می‌کنید:

\begin{latin}
\begin{lstlisting}
[me@linuxbox ~]$ printf "Line: %05d %15.3f Result: %+15d\n" 1071 3.14156295 32589
Line: 01071           3.142 Result:       +32589
\end{lstlisting}
\end{latin}

این مثال به‌خوبی تأثیر «حداقل عرض فیلد» را در ایجاد فواصل منظم میان بخش‌های مختلف خروجی نشان می‌دهد. در اینجا، عدد اول با صفرهای پیشین پر شده، عدد دوم تا سه رقم اعشار گرد شده و در فیلدی به عرض ۱۵ کاراکتر قرار گرفته، و عدد سوم نیز با حفظ علامت مثبت در یک فضای ۱۵ کاراکتری تراز شده است.

علاوه بر قالب‌بندی داده‌های عددی و متنی، از \cmd{printf} می‌توان برای تولید خودکار متن‌های ساختاریافته نیز بهره برد؛ برای نمونه، می‌توان یک رشته‌ی قالب شامل برچسب‌های زبان نشانه‌گذاری \en{HTML} تعریف کرد که در آن، به‌جای بخش‌های متغیر (مانند عنوان یا محتوای صفحه)، از مشخصه‌های تبدیل \cmd{\%s} استفاده شده باشد. با فراخوانیِ این دستور و ارسال مقادیر دلخواه به‌عنوان آرگومان، می‌توان یک فایل \en{HTML} ساده را به‌طور خودکار و در یک مرحله تولید کرد:

\begin{latin}
\begin{lstlisting}
[me@linuxbox ~]$ printf "<html>\n\t<head>\n\t\t<title>%s</title>\n
\t</head>\n\t<body>\n\t\t<p>%s</p>\n\t</body>\n</html>\n" "Page Tit
le" "Page Content"
<html>
	<head>
		<title>Page Title</title>
	</head>
	<body>
		<p>Page Content</p>
	</body>
</html>
\end{lstlisting}
\end{latin}

در این نمونه، استفاده از توالی‌های گریز \cmd{\textbackslash n} (خط جدید) و \cmd{\textbackslash t} (تورفتگی) در کنار مشخصه‌های تبدیل، امکان تولید خروجی‌های ساختاریافته و خوانا را فراهم کرده است. این رویکرد به‌ویژه در اسکریپت‌های توزیع‌شده که نیاز به تولید گزارش‌های \en{HTML} یا فایل‌های پیکربندی پویا دارند، بسیار سودمند است.

\section{سیستم‌های حروف‌چینی و آماده‌سازی سند}

تا این مرحله، ابزارهای مقدماتیِ قالب‌بندی متن را بررسی کردیم. این ابزارها برای پردازش‌های محدود و ساده کارآمدند، اما برای پروژه‌های بزرگ‌تر و پیچیده‌تر چه راهکاری وجود دارد؟ یکی از دلایل محبوبیت یونیکس در جوامع فنی و علمی، فراتر از ارائه‌ی محیطی قدرتمند، چندوظیفه‌ای و چندکاربره برای توسعه‌ی نرم‌افزار، بهره‌مندی از ابزارهایی بود که امکان تولید اسناد تخصصی، به‌ویژه انتشارات علمی و دانشگاهی، را فراهم می‌کردند. در واقع، طبق مستندات پروژه‌ی گنو، سیستم‌های آماده‌سازی سند نقشی کلیدی در تاریخچه‌ی توسعه‌ی یونیکس ایفا کرده‌اند.

نخستین نسخه‌ی یونیکس بر روی رایانه‌ی \en{PDP-7} در آزمایشگاه‌های بل توسعه یافت. در سال ۱۹۷۱، توسعه‌دهندگان برای توجیه هزینه‌ی خرید یک دستگاه جدید، یعنی \en{PDP-11}، پیشنهاد دادند که یک سیستم قالب‌بندی سند برای بخش ثبت اختراع (\en{Patent Department}) شرکت \en{AT\&T} پیاده‌سازی کنند. این برنامه‌ی اولیه، بازنویسی‌ای از \cmd{roff}، اثرِ داگلاس مک‌ایلروی (\en{Douglas McIlroy})، بود که توسط جوزف اف. اوسانا (\en{Joseph F. Ossanna}) نوشته شد. خودِ \cmd{roff} نیز ریشه در برنامه‌ای به نام \cmd{runoff} داشت که پیش‌تر توسط جروم ای. سالتزر (\en{Jerome E. Saltzer}) برای سیستم اشتراک زمانی سازگار (\en{CTSS}) در \en{MIT} نوشته شده بود؛ نامِ \cmd{roff} نیز برگرفته از همین واژه‌ی \cmd{runoff} است که در ادبیات آن دوران به معنای «تهیه‌ی یک نسخه‌ی چاپی» به کار می‌رفت.

در دنیای لینوکس و یونیکس، دو خانواده‌ی اصلی از سیستم‌های حروف‌چینی حضور دارند: نخست، ابزارهای مشتق‌شده از \cmd{roff} اصلی (مانند \cmd{nroff} و \cmd{troff})، و دوم، سیستم‌های مبتنی بر موتور حروف‌چینیِ دونالد کِنوث (\en{Donald Knuth}) با نام \en{TEX}. شایان ذکر است که سبک نگارشیِ نام \en{TEX} و تراز متفاوت حرف \en{E} در آن، عمدی و برگرفته از الفبای یونانی است.

برنامه‌ی \cmd{nroff} (مخفف \en{New Roff}) برای قالب‌بندی اسناد بر روی دستگاه‌هایی با قلم‌های تک‌فاصله (\en{Monospaced})، نظیر ترمینال‌های متنی و چاپگرهای سوزنی، طراحی شده بود؛ در دوران ظهور این ابزار، تقریباً تمامی چاپگرهای موجود در همین دسته قرار می‌گرفتند. نسخه‌ی پیشرفته‌تر، یعنی \cmd{troff} (مخفف \en{Typesetter Roff})، اسناد را برای خروجی روی دستگاه‌های حروف‌چینِ حرفه‌ای (\en{Phototypesetters}) آماده می‌کرد؛ دستگاه‌هایی که متونی با کیفیت چاپ تجاری تولید می‌کردند. امروزه اکثر چاپگرهای مدرن قادر به شبیه‌سازی خروجی این دستگاه‌ها هستند. خانواده‌ی \cmd{roff} همچنین ابزارهای جانبیِ متعددی نظیر \cmd{eqn} (برای فرمول‌نویسی ریاضی) و \cmd{tbl} (برای ترسیم جداول) را در بر می‌گیرد.

نسخه‌ی پایدار سیستم \en{TEX}، یعنی \en{TEX3}، نخستین بار در سال ۱۹۸۹ عرضه شد و تا حد زیادی جایگزین \cmd{troff} در پروژه‌های حروف‌چینیِ سنگین گردید. در این مبحث به \en{TEX} نخواهیم پرداخت؛ چرا که نه‌تنها ساختاری بسیار پیچیده دارد (به‌طوری که کتاب‌های متعددی درباره‌ی آن نگاشته شده است)، بلکه به‌صورت پیش‌فرض نیز در اکثر توزیع‌های مدرن لینوکس نصب نمی‌شود.

\begin{note}
\textbf{نکته:} کاربرانی که تمایل به نصب و استفاده از \en{TEX} دارند، می‌توانند مجموعه‌ی \cmd{texlive} را، که در اکثر مخازن رسمی توزیع‌ها موجود است، به همراه واژه‌پرداز گرافیکی \en{LyX} بررسی کنند.
\end{note}

\subsection{دستور groff: موتور حروف‌چینی گنو}

\cmd{groff} مجموعه‌ای از نرم‌افزارهاست که پیاده‌سازی پروژه‌ی گنو (\en{GNU}) از سیستم \cmd{troff} را ارائه می‌دهد. این مجموعه همچنین شامل اسکریپت‌هایی برای شبیه‌سازی \cmd{nroff} و سایر اعضای خانواده‌ی بزرگ \cmd{roff} است.

هرچند \cmd{roff} و نسخه‌های مشتق شده از آن برای تولید اسناد قالب‌بندی‌شده طراحی شده‌اند، اما متدولوژی آن‌ها با ذهنیت کاربران امروزی تفاوت دارد. امروزه اکثر اسناد با واژه‌پردازهای گرافیکی تولید می‌شوند که فرآیند نگارش و صفحه‌بندی را به‌صورت هم‌زمان (\en{WYSIWYG - What You See Is What You Get}) مدیریت می‌کنند. اما پیش از ظهور این ابزارها، اسناد طی یک فرآیند دو مرحله‌ای تولید می‌شدند: ابتدا متن خام در یک ویرایشگر متن نگارش می‌شد و سپس یک پردازشگر مانند \cmd{troff} وظیفه‌ی اعمال قالب‌بندی را بر عهده می‌گرفت. در این روش، دستورالعمل‌های قالب‌بندی از طریق یک «زبان نشانه‌گذاری» (\en{Markup Language}) در میان متن اصلی گنجانده می‌شدند.

نزدیک‌ترین معادل مدرن برای این فرآیند، صفحات وب هستند؛ جایی که متن در یک ویرایشگر نوشته شده و سپس مرورگر وب با استفاده از زبان نشانه‌گذاری \en{HTML}، آن را به فرمت بصری نهایی تبدیل و نمایش می‌دهد.

هدف ما در اینجا آموزش جامع \cmd{groff} نیست، چرا که بسیاری از دستورات زبان نشانه‌گذاری آن به جزئیات بسیار تخصصی در صنعت حروف‌چینی مربوط می‌شود. در مقابل، تمرکز ما بر یکی از بسته‌های ماکرو (\en{Macro Packages}) خواهد بود که همچنان کاربرد گسترده‌ای دارد. این بسته‌ها، دستورات سطح پایین و پیچیده‌ی \cmd{groff} را در قالب مجموعه‌ای از فرمان‌های سطح بالا و ساده‌تر انتزاع می‌کنند که استفاده از این سیستم را بسیار تسهیل می‌بخشد.

بیایید نگاهی به ساختار درونی یک صفحه‌ی راهنما (\en{Man Page}) بیندازیم. این مستندات در مسیر \file{/usr/share/man/} به‌صورت فایل‌های متنی فشرده با فرمت \cmd{gzip} نگهداری می‌شوند. با بررسی محتوای یکی از این فایل‌ها (برای مثال، صفحه‌ی راهنمای دستور \cmd{ls} در بخش ۱)، با کدهای زیر مواجه می‌شویم:

\begin{latin}
\begin{lstlisting}
[me@linuxbox ~]$ zcat /usr/share/man/man1/ls.1.gz | head
.\" DO NOT MODIFY THIS FILE! It was generated by help2man 1.47.3.
.TH LS "1" "January 2025" "GNU coreutils 8.28" "User Commands"
.SH NAME
ls \- list directory contents
.SH SYNOPSIS
.B ls
[\fI\,OPTION\/fR]... [\fI\,FILE\/\fR]...
.SH DESCRIPTION
.\" Add any additional description here
.PP
\end{lstlisting}
\end{latin}

با مقایسه‌ی خروجی استاندارد یک صفحه‌ی راهنما با کد منبع آن، می‌توان به‌وضوح ارتباط میان زبان نشانه‌گذاری و نتیجه‌ی نهایی را درک کرد:

\begin{latin}
\begin{lstlisting}
[me@linuxbox ~]$ man ls | head
LS(1)                    User Commands                          LS(1)

NAME
       ls - list directory contents

SYNOPSIS
       ls [OPTION]... [FILE]...
\end{lstlisting}
\end{latin}

علت اهمیت این ساختار آن است که تمامی صفحات راهنمای لینوکس توسط موتور \cmd{groff} و با بهره‌گیری از بسته‌ی ماکروی تخصصی \cmd{mandoc} پردازش می‌شوند. در واقع، می‌توان عملکرد داخلی دستور \cmd{man} را با استفاده از توالی فرمان‌های زیر شبیه‌سازی کرد:

\begin{latin}
\begin{lstlisting}
[me@linuxbox ~]$ zcat /usr/share/man/man1/ls.1.gz | groff -mandoc -Tascii | head
LS(1)                     User Commands                          LS(1)

NAME
       ls - list directory contents

SYNOPSIS
       ls [OPTION]... [FILE]
\end{lstlisting}
\end{latin}

در این دستور، ابزار \cmd{groff} با گزینه \cmd{-mandoc} فراخوانی شده تا مجموعه‌ ماکروهای لازم برای تفسیر مستندات سیستم بارگذاری شوند. همچنین گزینه \cmd{-Tascii} درایور خروجی را بر روی فرمت متنی \en{ASCII} تنظیم می‌کند تا نتیجه به‌صورت متنی در ترمینال قابل نمایش باشد.

ابزار \cmd{groff} قادر است خروجی خود را در چندین قالب مختلف تولید کند. در صورتی که هیچ فرمت خاصی توسط کاربر تعیین نشود، این برنامه به‌طور پیش‌فرض خروجی را در قالب \en{PostScript} تولید می‌کند:

\begin{latin}
\begin{lstlisting}
[me@linuxbox ~]$ zcat /usr/share/man/man1/ls.1.gz | groff -mandoc | head
%!PS-Adobe-3.0
%%Creator: groff version 1.18.1
%%CreationDate: Thu Feb  5 13:44:37 2025
%%DocumentNeededResources: font Times-Roman
%%+ font Times-Bold
%%+ font Times-Italic
%%DocumentSuppliedResources: procset grops 1.18 1
%%Pages: 4
%%PageOrder: Ascend
%%Orientation: Portrait
\end{lstlisting}
\end{latin}

در فصل گذشته به‌طور اجمالی درباره‌ی \en{PostScript} صحبت کردیم و در فصل آتی نیز مجدداً به آن خواهیم پرداخت. \en{PostScript} یک زبان توصیف صفحه (\en{Page Description Language}) است که به‌منظور تبیین دقیق محتویات یک صفحه‌ی چاپی برای دستگاه‌های حروف‌چین و چاپگرهای پیشرفته طراحی شده است.

چنانچه خروجی دستور فوق را در یک فایل ذخیره نمایید (با این فرض که از یک محیط دسکتاپ گرافیکی استفاده می‌کنید)، آیکون فایل مربوطه باید بر روی دسکتاپ ظاهر شود:

\begin{latin}
\begin{lstlisting}
[me@linuxbox ~]$ zcat /usr/share/man/man1/ls.1.gz | groff -mandoc > ~/Desktop/ls.ps
\end{lstlisting}
\end{latin}

با دوبار کلیک بر روی این آیکون، برنامه‌ی نمایشگر اسناد (\en{Document Viewer}) اجرا شده و فایل را در حالت رندرشده‌ی نهایی به نمایش می‌گذارد؛ درست همان‌گونه که در شکل ۴ مشاهده می‌شود.

\bookfigure{img-04.png}{رندرینگ خروجی \en{PostScript} توسط یک نمایشگر اسناد در محیط گرافیکی \en{GNOME}}

خروجی حاصل، یک نسخه‌ی حروف‌چینی شده و آراسته از صفحه‌ی راهنمای دستور \cmd{ls} است. برای سهولت بیشتر، می‌توان فایل \en{PostScript} تولید شده را با استفاده از دستور زیر به قالب \en{PDF} (\en{Portable Document Format}) تبدیل کرد:

\begin{latin}
\begin{lstlisting}
[me@linuxbox ~]$ ps2pdf ~/Desktop/foo.ps ~/Desktop/ls.pdf
\end{lstlisting}
\end{latin}

ابزار \cmd{ps2pdf} بخشی از مجموعه‌ی نرم‌افزاری \en{Ghostscript} است که در اکثر توزیع‌های لینوکس مجهز به سیستم چاپ، به‌صورت پیش‌فرض نصب می‌شود.

\begin{note}
\textbf{نکته:} سیستم‌عامل لینوکس شامل ابزارهای خط فرمان متعددی جهت تبدیل قالب‌های فایلی به یکدیگر است. این برنامه‌ها معمولاً از الگوی نام‌گذاری \cmd{format2format} پیروی می‌کنند. برای مشاهده‌ی فهرستی از این ابزارها، می‌توانید دستور زیر را اجرا کنید:
\begin{latin}
\begin{lstlisting}
ls /usr/bin/*[[:alpha:]]2[[:alpha:]]*
\end{lstlisting}
\end{latin}
همچنین جست‌وجو برای برنامه‌هایی که با الگوی نام‌گذاری \cmd{formattoformat} تعریف شده‌اند نیز می‌تواند سودمند باشد.
\end{note}

در آخرین تمرین با \cmd{groff} به فایل آشنای \file{distros.txt} بازمی‌گردیم. این بار از برنامه‌ی \cmd{tbl} که ابزاری تخصصی برای قالب‌بندی جداول در خروجی‌های حروف‌چینی است، جهت آرایش فهرست توزیع‌های لینوکس بهره خواهیم برد. برای این منظور، از اسکریپت \cmd{sed} که پیش‌تر طراحی کردیم استفاده می‌کنیم تا نشانه‌گذاری‌های (\en{Markup}) لازم را به جریان متن افزوده و خروجی را جهت پردازش نهایی به \cmd{groff} هدایت کنیم.

نخست باید اسکریپت \cmd{sed} خود را به‌گونه‌ای بازنویسی کنیم که المان‌های نشانه‌گذاری لازم (که در ادبیات \cmd{groff} با عنوان «دستور» یا \en{Request} شناخته می‌شوند) را به متن تزریق کند. با استفاده از یک ویرایشگر متن، فایل \file{distros.sed} را مطابق الگوی زیر اصلاح می‌کنیم:

\begin{latin}
\begin{lstlisting}
# sed script to produce Linux distributions report

1 i\
.TS\
center box;\
cb s s\
cb cb cb\
l n c.\
Linux Distributions Report\
=\
Name    Version    Released\
-\
s/\([0-9]\{2\}\)\/\([0-9]\{2\}\)\/\([0-9]\{4\}\)$/\3-\1-\2/
$ a\
.TE
\end{lstlisting}
\end{latin}

شایان ذکر است که برای عملکرد صحیح اسکریپت، فیلدهای \en{Name}، \en{Version} و \en{Released} الزاماً باید با کاراکتر \en{Tab} از یکدیگر جدا شوند. این نسخه از اسکریپت را با نام \file{distros-tbl.sed} ذخیره می‌کنیم. پیش‌پردازنده‌ی \cmd{tbl} از درخواست‌های \cmd{.TS} (مخفف \en{Table Start}) و \cmd{.TE} (مخفف \en{Table End}) برای تعیین محدوده جدول استفاده می‌کند.

خطوط بلافاصله پس از دستور \cmd{.TS} پیکربندی عمومی جدول را تعیین می‌کنند؛ در این مثال، جدول در مرکز صفحه قرار گرفته و محصور در یک کادر (\en{box}) ترسیم می‌شود. سطر بعدی چیدمان ستون‌ها را تعریف کرده و سایر خطوط نیز ساختار هر سطر را توصیف می‌کنند. اکنون با اجرای مجدد پایپ‌لاین (\en{Pipeline}) تولید گزارش به‌وسیله‌ی اسکریپت جدید، خروجی به شکل زیر ظاهر خواهد شد:

\begin{latin}
\begin{lstlisting}
[me@linuxbox ~]$ sort -k 1,1 -k 2n distros.txt | sed -f distros-tbl.sed | groff -t -T ascii
               +---------------------------------------+
               | Linux Distributions Report            |
               +---------------------------------------+
               | Name       Version     Released       |
               +---------------------------------------+
               | Fedora    5           2006-03-20      |
               | Fedora    6           2006-10-24      |
               | Fedora    7           2007-05-31      |
               | Fedora    8           2007-11-08      |
               | Fedora    9           2008-05-13      |
               | Fedora    10          2008-11-25      |
               | SUSE      10.1        2006-05-11      |
               | SUSE      10.2        2006-12-07      |
               | SUSE      10.3        2007-10-04      |
               | SUSE      11.0        2008-06-19      |
               | Ubuntu    6.06        2006-06-01      |
               | Ubuntu    6.10        2006-10-26      |
               | Ubuntu    7.04        2007-04-19      |
               | Ubuntu    7.10        2007-10-18      |
               | Ubuntu    8.04        2008-04-24      |
               | Ubuntu    8.10        2008-10-30      |
               +---------------------------------------+
\end{lstlisting}
\end{latin}

افزودن گزینه \cmd{-t} به \cmd{groff} به این برنامه دستور می‌دهد که جریان متن ورودی را پیش از پردازش نهایی، از فیلتر \cmd{tbl} عبور دهد. همچنین، گزینه‌ی \cmd{-T ascii} برای تنظیم قالب خروجی بر روی استاندارد \en{ASCII} (به‌جای قالب پیش‌فرض \en{PostScript}) استفاده شده است.

خروجی متنی بالا، حداکثر کیفیتی است که می‌توان در محیط یک ترمینال متنی یا چاپگرهای تک‌فاصله (\en{Monospaced}) به آن دست یافت. با این حال، اگر خروجی را بر روی قالب \en{PostScript} تنظیم کرده و آن را در یک محیط گرافیکی مشاهده کنیم، به نتیجه‌ای بسیار حرفه‌ای‌تر و آراسته‌تر دست خواهیم یافت؛ همان‌گونه که در شکل ۵ نمایش داده شده است:

\begin{latin}
\begin{lstlisting}
[me@linuxbox ~]$ sort -k 1,1 -k 2n distros.txt | sed -f distros-tbl.sed | groff -t > ~/Desktop/foo.ps
\end{lstlisting}
\end{latin}

\bookfigure{img-05.png}{نمای نهایی جدول پردازش شده در یک نمایشگر سند}

در این حالت، \cmd{groff} از فونت‌های متناسب (\en{Proportional Fonts}) و خطوط ترسیمی دقیق برای کادربندی استفاده می‌کند که تفاوت فاحشی با خروجی‌های متنی ساده دارد.

\section{جمع‌بندی}

با توجه به جایگاه محوری متن در ساختار سیستم‌عامل‌های شبه‌یونیکس، وجود ابزارهای متنوع برای پردازش و آرایش متون کاملاً منطقی به نظر می‌رسد. همان‌طور که بررسی شد، اکوسیستم لینوکس طیف گسترده‌ای از این قابلیت‌ها را ارائه می‌دهد؛ ابزارهای ساده‌ای نظیر \cmd{fmt} و \cmd{pr} برای تولید اسناد کوتاه و سریع در اسکریپت‌نویسی بسیار کارآمد هستند، در حالی که مجموعه‌ی \cmd{groff} و ابزارهای مکمل آن، توانایی مدیریت پروژه‌های سنگین نظیر نگارش کتاب‌های کامل را دارا می‌باشند. اگرچه ممکن است نگارش مقالات فنی با استفاده از ابزارهای خط فرمان برای همه‌ی کاربران مرسوم نباشد (هرچند متخصصان بسیاری این روش را ترجیح می‌دهند)، اما آگاهی از این توانمندی‌ها و تسلط بر آن‌ها، ابزار قدرتمندی را در اختیار مدیران سیستم و توسعه‌دهندگان قرار می‌دهد.

\section{منابع مطالعه تکمیلی}

راهنمای کاربری \cmd{groff}:

\begin{latin}
\url{http://www.gnu.org/software/groff/manual/}
\end{latin}

مقاله‌نویسی با \cmd{nroff} و مجموعه ماکروی \en{me}:

\begin{latin}
\url{http://docs.freebsd.org/44doc/usd/19.memacros/paper.pdf}
\end{latin}

راهنمای مرجع مجموعه ماکروی \en{me}:

\begin{latin}
\url{http://docs.freebsd.org/44doc/usd/20.meref/paper.pdf}
\end{latin}

همچنین، مطالعه مقالات زیر از ویکی‌پدیا توصیه می‌شود:

\begin{latin}
\url{http://en.wikipedia.org/wiki/TeX}\\
\url{http://en.wikipedia.org/wiki/Donald_Knuth}\\
\url{http://en.wikipedia.org/wiki/Typesetting}
\end{latin}